\documentclass[11pt,a4paper]{article}
\usepackage[utf8]{inputenc}
\usepackage{amsmath,amssymb,amsthm}
\usepackage{geometry}
\usepackage{booktabs}
\usepackage{hyperref}
\usepackage{microtype}
\usepackage{xcolor}

\geometry{margin=1in}

\title{\textbf{PartitionViz: Computational and Interactive Investigations in the Theory of Partitions}\\
\large Independent Project (IP) Milestone Report --- Chapters 1 \& 2}
\author{\textbf{Jatin Pawar}\\
Department of Mathematics / Computer Science\\
Under the Supervision of Course Professor / Advisor}
\date{Academic Term 2026}

\newtheorem{theorem}{Theorem}[section]
\newtheorem{definition}[theorem]{Definition}
\newtheorem{corollary}[theorem]{Corollary}
\newtheorem{example}[theorem]{Example}

\begin{document}

\maketitle

\begin{abstract}
We present \textbf{PartitionViz}, an interactive computational mathematics framework designed to investigate the combinatorial bijections, geometric involutions, and generating function identities presented in George E. Andrews's \emph{The Theory of Partitions} (Cambridge University Press). In this milestone, we have developed three integrated research modules: (1) Ferrers diagram transposition and the bounded partition duality of Andrews's Theorem 1.4, including Sylvester's hook unfolding for self-conjugate partitions; (2) Franklin's sign-reversing involution (1881) demonstrating exact parity cancellation in Euler's pentagonal number series $\prod_{k=1}^\infty (1-q^k)$ and powering Euler's $O(N\sqrt{N})$ recurrence for $p(n)$; and (3) Series-product identities from Chapter 2, including formal power series arithmetic in $\mathbb{Z}[[q]]$, Euler's odd-distinct algebraic identity, and Andrews's Durfee square series expansion $\sum_{d=0}^\infty \frac{q^{d^2}}{(q; q)_d^2} = \sum p(n) q^n$. The system features an independent mathematical engine verified by 50 automated unit tests, interactive visualizations with step-by-step geometric transformations, and defensible computational findings.
\end{abstract}

\section{Introduction and Theoretical Scope}
The theory of integer partitions, initiated systematically by Leonhard Euler in 1748, is foundational to modern combinatorics, representation theory, and modular forms \cite{andrews1976}. A \textbf{partition} of a positive integer $n$ is a finite non-increasing sequence of positive integers:
\begin{equation}
\lambda = (\lambda_1, \lambda_2, \dots, \lambda_r), \quad \lambda_1 \ge \lambda_2 \ge \dots \ge \lambda_r > 0, \quad \sum_{i=1}^r \lambda_i = n.
\end{equation}
The number of partitions of $n$ is denoted by $p(n)$, with generating function:
\begin{equation}
\sum_{n=0}^\infty p(n) q^n = \prod_{k=1}^\infty \frac{1}{1 - q^k}.
\end{equation}
Rather than confining our Independent Project to passive theoretical review, we have established a computational laboratory adhering to the cycle:
\begin{equation*}
\text{Theorem} \longrightarrow \text{Algorithmic Engine} \longrightarrow \text{Computational Experiment} \longrightarrow \text{Deduction}.
\end{equation}

\section{Module 1: Ferrers Graphs and Conjugation Duality}

\subsection{Mathematical Formulation}
For a partition $\lambda$, the \textbf{Ferrers diagram} represents $\lambda$ as an array of dots with $\lambda_i$ dots in row $i$. The \textbf{conjugate partition} $\lambda' = (\lambda'_1, \dots, \lambda'_m)$ is defined by:
\begin{equation}
\lambda'_i = \#\{j : \lambda_j \ge i\}.
\end{equation}
Geometrically, $\lambda'$ is obtained by reflecting the Ferrers diagram across its main diagonal.

\subsection{Andrews's Theorem 1.4}
\begin{theorem}[Andrews, Theorem 1.4, p.~8]
The number of partitions of $n$ into at most $m$ parts equals the number of partitions of $n$ whose largest part is at most $m$.
\end{theorem}
\begin{proof}[Combinatorial Mechanism]
The conjugation map $\phi_{\text{conj}} : \lambda \mapsto \lambda'$ provides an explicit bijection. If $\ell(\lambda) \le m$, then the number of rows of $\lambda$ is at most $m$, which implies the largest column height of $\lambda'$ (and hence $\lambda'_1$) is at most $m$. Because $(\lambda')' = \lambda$, the map is an exact involution on $\mathcal{P}(n)$.
\end{proof}

\subsection{Empirical Deductions}
\begin{enumerate}
    \item \textbf{Sylvester's Fixed-Point Decay:} By Sylvester's theorem (Andrews Ch.~1, Ex.~8), the self-conjugate partitions ($\lambda = \lambda'$) are equinumerous with partitions into distinct odd parts. While $p(n) \sim \frac{1}{4n\sqrt{3}} e^{\pi \sqrt{2n/3}}$, self-conjugate partitions grow only as $e^{\pi \sqrt{n/3}}$. The proportion $sc(n)/p(n)$ decays to zero exponentially fast at rate $e^{-\pi(\sqrt{2}-1)\sqrt{n/3}}$.
    \item \textbf{Durfee Square Scaling:} The Durfee square side $d(\lambda) = \max\{i : \lambda_i \ge i\}$ satisfies $d(\lambda) = d(\lambda')$. Computationally, its expected value scales as $E[d(n)] \approx \frac{\sqrt{6}\ln 2}{\pi}\sqrt{n} \approx 0.5407\sqrt{n}$, reflecting the limiting shape of random partitions.
    \item \textbf{Durfee Partition Decomposition:} Every partition decomposes uniquely into a $d \times d$ square, a partition $\lambda_R$ of dots to the right (at most $d$ parts), and a partition $\lambda_B$ of dots below (parts at most $d$), establishing:
    \begin{equation}
    |\lambda| = d^2 + |\lambda_R| + |\lambda_B|.
    \end{equation}
\end{enumerate}

\section{Module 2: Franklin's Involution and Euler's Pentagonal Series}

\subsection{The Analytic Identity}
In Chapter 1, Section 1.3, Andrews introduces Euler's celebrated pentagonal number theorem:
\begin{equation}
\prod_{k=1}^\infty (1 - q^k) = 1 + \sum_{k=1}^\infty (-1)^k \left[ q^{k(3k-1)/2} + q^{k(3k+1)/2} \right].
\end{equation}
Combinatorially, the coefficient of $q^n$ is $p_e(D, n) - p_o(D, n)$, where $p_e$ (resp. $p_o$) counts partitions of $n$ into an even (resp. odd) number of distinct parts.

\subsection{Franklin's Sign-Reversing Involution (1881)}
Let $\lambda \in D(n)$. Define the base $b = \lambda_r$ and the slope $s = \max\{s : \lambda_i = \lambda_1 - (i-1) \text{ for } 1 \le i \le s\}$.
\begin{itemize}
    \item If $b \le s$ (and not exceptional): delete part $b$ and add $1$ to the first $b$ parts ($\ell \to \ell - 1$).
    \item If $b > s$ (and not exceptional): subtract $1$ from the first $s$ parts and append a new part $s$ ($\ell \to \ell + 1$).
    \item \textbf{Exceptional Cases ($s = r$):} When the base meets the slope, the move fails if $b = s$ (giving $n = \omega(k)$) or $b = s + 1$ (giving $n = \omega(-k)$).
\end{itemize}

\subsection{Euler's Recurrence for $p(n)$ and Complexity Analysis}
Equating $(\sum p(n) q^n) \prod (1 - q^k) = 1$ gives Andrews Corollary 1.8:
\begin{equation}
p(n) = \sum_{k \ge 1} (-1)^{k+1} \left[ p(n - \omega(k)) + p(n - \omega(-k)) \right].
\end{equation}
\begin{itemize}
    \item \textbf{Single Evaluation:} Since $\omega(k) \approx \frac{3}{2}k^2 \le n$, only $O(\sqrt{n})$ terms ($\approx 2\sqrt{2n/3}$) are required.
    \item \textbf{Sequential Table:} Evaluating $p(0), \dots, p(N)$ takes $\sum_{n=1}^N O(\sqrt{n}) = O(N^{3/2}) = O(N\sqrt{N})$ operations.
    \item \textbf{Naive Enumeration:} In contrast, generating partitions requires exponential time $O(e^{\pi\sqrt{2N/3}})$.
\end{itemize}

\begin{table}[h]
\centering
\caption{Exact Recurrence Verification for Select Values of $p(n)$}
\begin{tabular}{rrrrc}
\toprule
$n$ & Naive Search Space & Exact $p(n)$ & Pentagonal Terms Used & Computation Time \\
\midrule
10 & 42 & 42 & 5 & $< 0.05$ ms \\
20 & 627 & 627 & 7 & $< 0.05$ ms \\
50 & 204,226 & 204,226 & 11 & $< 0.10$ ms \\
100 & 190,569,292 & 190,569,292 & 15 & $< 0.20$ ms \\
\bottomrule
\end{tabular}
\end{table}

\section{Module 3: Series-Product Identities (Andrews Chapter 2)}

\subsection{Formal Power Series in $\mathbb{Z}[[q]]$}
In Chapter 2, Andrews formalizes generating functions as formal power series. Truncating an infinite product $\prod_{k=1}^\infty (1 \pm q^k)^{\pm 1}$ at $M \ge N$ factors stabilizes the exact integer coefficients up to degree $N$:
\begin{equation}
\prod_{k=1}^\infty \frac{1}{1 - q^k} \equiv \prod_{k=1}^N \frac{1}{1 - q^k} \pmod{q^{N+1}}.
\end{equation}

\subsection{Euler's Odd-Distinct Identity}
\begin{theorem}[Euler; Andrews Corollary 1.2, p.~5]
The number of partitions of $n$ into distinct parts equals the number of partitions of $n$ into odd parts.
\end{theorem}
The algebraic proof follows by factorizing $1 + q^k = \frac{1 - q^{2k}}{1 - q^k}$, leading to telescoping cancellation of all even factors:
\begin{equation}
\prod_{k=1}^\infty (1 + q^k) = \prod_{k=1}^\infty \frac{1 - q^{2k}}{1 - q^k} = \prod_{k=1}^\infty \frac{1}{1 - q^{2k-1}}.
\end{equation}

\subsection{Durfee Square Series Expansion}
\begin{theorem}[Andrews Section 2.2, eq.~2.2.9, p.~28]
The unrestricted partition generating function expands as a sum of rational functions:
\begin{equation}
\sum_{n=0}^\infty p(n) q^n = \sum_{d=0}^\infty \frac{q^{d^2}}{(1-q)^2(1-q^2)^2\dots(1-q^d)^2}.
\end{equation}
\end{theorem}
Combinatorially, every partition has a Durfee square of side $d$ ($q^{d^2}$), a right partition of $\le d$ parts ($\frac{1}{(q; q)_d}$), and a bottom partition with parts $\le d$ (also $\frac{1}{(q; q)_d}$ by conjugation).

\section{Software Architecture and Quality Assurance}
PartitionViz is designed with a strict separation between presentation and computation:
\begin{itemize}
    \item \textbf{\texttt{partition\_core}}: Independent Python package containing enumeration algorithms, exact integer arithmetic, involution logic, formal power series arithmetic, and asymptotic estimators.
    \item \textbf{\texttt{ui}}: Modular Streamlit presentation layer featuring an academic research aesthetic, Plotly Ferrers renderers, native LaTeX equation rendering, and step-by-step animations.
    \item \textbf{One-Command Execution}: The entire suite runs via \texttt{streamlit run app.py}.
    \item \textbf{Test Suite}: 50 automated tests in \texttt{pytest} verifying sum preservation, involution symmetry, distinctness, series-product coefficients, and Durfee square expansions.
\end{itemize}

\begin{thebibliography}{9}
\bibitem{andrews1976}
G.~E.~Andrews, \emph{The Theory of Partitions}, Encyclopedia of Mathematics and Its Applications, Vol.~2, Addison-Wesley, 1976 (Reissued by Cambridge University Press, 1998).
\bibitem{franklin1881}
F.~Franklin, \emph{Sur le d\'eveloppement du produit infini $(1-x)(1-x^2)(1-x^3)\dots$}, Comptes Rendus Acad. Sci. Paris, 92 (1881), 448--450.
\bibitem{sylvester1882}
J.~J.~Sylvester, \emph{A Constructive Theory of Partitions}, Amer. J. Math., 5 (1882), 251--330.
\end{thebibliography}

\end{document}
